% FAC-EN-U0035
% French source: fac_body.tex lines 1321--1343.
% Authority: PDF physical pp. 35--36; printed pp. 230--231 (no. 37 complete).
% Transcription repair: FAC-FR-T0023. Printed-source repair: FAC-FR-C0027.

\phantomsection\label{fac:II:2}
\section*{\S2. Coherent algebraic sheaves}

\phantomsection\label{fac:II:2:no:37}
\textbf{37. The sheaf of local rings of an algebraic variety.}
Return to the notation of \hyperref[fac:II:1:no:31]{no.~31}: let $X=K^r$, and let $\mathcal{O}$ be the sheaf of local rings of $X$.

\phantomsection\label{fac:II:2:no:37:lemma:1}
\textsc{Lemma 1.} \textit{The sheaf $\mathcal{O}$ is a coherent sheaf of rings, in the sense of \hyperref[fac:I:2:no:15]{no.~15}.}

Let $x\in X$, let $U$ be a neighborhood of $x$, and let $f_1,\ldots,f_p$ be sections of $\mathcal{O}$ over $U$, that is, rational functions regular at every point of $U$. We must show that the sheaf of relations among $f_1,\ldots,f_p$ is a sheaf of finite type over $\mathcal{O}$.

After replacing $U$ by a smaller neighborhood, we may suppose that
\[
f_i=P_i/Q,
\]
where the $P_i$ and $Q$ are polynomials and $Q$ does not vanish on $U$. Now let $y\in U$, and let $g_i\in\mathcal{O}_y$ be such that
\[
\sum_{i=1}^{i=p}g_i f_i=0
\]
in a neighborhood of $y$. We may write $g_i=R_i/S$, where the $R_i$ and $S$ are polynomials and $S$ does not vanish at $y$. The relation ``$\sum_{i=1}^{i=p}g_i f_i=0$ in a neighborhood of $y$'' is equivalent to the relation ``$\sum_{i=1}^{i=p}R_iP_i=0$ in a neighborhood of $y$,'' which in turn is equivalent to
\[
\sum_{i=1}^{i=p}R_iP_i=0.
\]
Since the module of relations among the polynomials $P_i$ is of finite type, because the polynomial ring is Noetherian, it follows that the sheaf of relations among the $f_i$ is of finite type.

Now let $V$ be a closed subvariety of $X=K^r$. For every $x\in X$, let $\mathcal{J}_x(V)$ be the ideal of $\mathcal{O}_x$ consisting of the elements $f\in\mathcal{O}_x$ whose restriction to $V$ is zero in a neighborhood of $x$; thus $\mathcal{J}_x(V)=\mathcal{O}_x$ if $x\notin V$. The $\mathcal{J}_x(V)$ form a subsheaf $\mathcal{J}(V)$ of $\mathcal{O}$.

\phantomsection\label{fac:II:2:no:37:lemma:2}
\textsc{Lemma 2.} \textit{The sheaf $\mathcal{J}(V)$ is a coherent sheaf of $\mathcal{O}$-modules.}

Let $I(V)$ be the ideal of $K[X_1,\ldots,X_r]$ consisting of the polynomials $P$ which vanish on $V$. By \hyperref[fac:II:1:no:31:prop:4]{Proposition~4 of no.~31},
\[
\mathcal{J}_x(V)=I(V)\mathcal{O}_x
\]
for every $x\in V$, and this formula also holds for $x\notin V$, as is immediately seen. Since the ideal $I(V)$ is generated by finitely many elements, the sheaf $\mathcal{J}(V)$ is of finite type and hence coherent by \hyperref[fac:II:2:no:37:lemma:1]{Lemma~1} and \hyperref[fac:I:2:no:15:prop:8]{Proposition~8 of no.~15}.

We shall now extend \hyperref[fac:II:2:no:37:lemma:1]{Lemma~1} to an arbitrary algebraic variety.

\phantomsection\label{fac:II:2:no:37:prop:1}
\textsc{Proposition 1.} \textit{If $V$ is an algebraic variety, then the sheaf $\mathcal{O}_V$ is a coherent sheaf of rings on $V$.}

The question is local, so we may suppose that $V$ is a closed subvariety of affine space $K^r$. By \hyperref[fac:II:2:no:37:lemma:2]{Lemma~2}, $\mathcal{J}(V)$ is a coherent sheaf of ideals. Hence $\mathcal{O}/\mathcal{J}(V)$ is a coherent sheaf of rings on $X$, by \hyperref[fac:I:2:no:16:thm:3]{Theorem~3 of no.~16}. This sheaf of rings is zero outside $V$, and its restriction to $V$ is precisely $\mathcal{O}_V$ (\hyperref[fac:II:1:no:31]{no.~31}). It follows that $\mathcal{O}_V$ is a coherent sheaf of rings on $V$ (\hyperref[fac:I:2:no:17:prop:11:cor]{no.~17, corollary to Proposition~11}).

\phantomsection\label{fac:II:2:no:37:remark}
\textsc{Remark.} It is clear that \hyperref[fac:II:2:no:37:prop:1]{Proposition~1} holds, more generally, for every prealgebraic variety.
